%==============================================================================
% Chapitre 33 - Nutation
%==============================================================================

\section{Introduction}

Dans le chapitre précédent, nous avons étudié la précession.

Nous avons vu qu'un projectile en rotation ne conserve pas exactement une
orientation fixe mais que son axe effectue un mouvement lent autour de la
direction moyenne du vol.

Cependant, un projectile réel présente généralement un second mouvement
oscillatoire.

Ce mouvement est appelé nutation.

\begin{aretenir}

La nutation est une oscillation superposée au mouvement de précession.

\end{aretenir}

\begin{figure}[htbp]
\centering
\begin{tikzpicture}

\draw[dashed] (0,-2) -- (0,2);

\draw[thick] (0,0) -- (2,1.5);

\draw[dashed] (0,0) circle (1);

\draw[->] (1,0) arc (0:60:1);

\node at (2.5,1.7) {Axe du projectile};

\end{tikzpicture}
\caption{Principe de la nutation.}
\label{fig:nutation}
\end{figure}

\section{Origine du terme}

Le mot nutation provient du latin \emph{nutare}, signifiant « osciller » ou
« hocher la tête ».

Le terme est utilisé dans plusieurs domaines :

\begin{itemize}
    \item mécanique ;
    \item astronomie ;
    \item dynamique des gyroscopes ;
    \item balistique.
\end{itemize}

Dans tous les cas, il désigne une oscillation secondaire venant se superposer
à un mouvement principal.

\section{Précession et nutation}

La précession et la nutation sont souvent observées simultanément.

La précession constitue généralement :

\begin{itemize}
    \item un mouvement lent ;
    \item relativement régulier.
\end{itemize}

La nutation correspond plutôt :

\begin{itemize}
    \item à une oscillation plus rapide ;
    \item à une amplitude plus faible.
\end{itemize}

\section{Exemple de la toupie}

Une toupie constitue à nouveau un excellent exemple.

Lorsqu'elle tourne :

\begin{itemize}
    \item son axe précessionne ;
    \item de petites oscillations apparaissent également.
\end{itemize}

Ces oscillations correspondent à la nutation.

\section{Apparition lors du tir}

La nutation apparaît principalement parce que le projectile ne quitte jamais
le canon dans un état parfaitement idéal.

À la sortie de bouche, plusieurs perturbations peuvent exister :

\begin{itemize}
    \item vibrations du canon ;
    \item dissymétries aérodynamiques ;
    \item perturbations des gaz ;
    \item défauts géométriques.
\end{itemize}

Ces perturbations initiales alimentent les oscillations.

\section{Mouvement de l'axe}

L'axe du projectile n'est pas parfaitement immobile.

Il oscille autour de sa position moyenne.

Cette oscillation se superpose au mouvement de précession.

La trajectoire réelle de l'axe devient alors plus complexe.

\section{Angle de nutation}

La nutation est souvent décrite à l'aide d'un angle.

Cet angle caractérise l'écart entre :

\begin{itemize}
    \item l'orientation instantanée ;
    \item l'orientation moyenne.
\end{itemize}

Dans la plupart des cas, cet angle reste faible.

\section{Amortissement aérodynamique}

L'air tend généralement à amortir la nutation.

Au cours du vol :

\begin{itemize}
    \item l'amplitude diminue ;
    \item le mouvement devient plus régulier ;
    \item le projectile se stabilise progressivement.
\end{itemize}

Ce phénomène est parfois appelé amortissement dynamique.

\section{Distance de mise en sommeil}

Après la sortie du canon, le projectile traverse souvent une phase de
stabilisation.

Pendant cette période :

\begin{itemize}
    \item la précession est importante ;
    \item la nutation est importante ;
    \item les oscillations s'atténuent progressivement.
\end{itemize}

On parle parfois de mise en sommeil (\emph{going to sleep}).

\begin{aretenir}

Un projectile ne vole généralement pas dans son état le plus stable dès la
sortie du canon.

\end{aretenir}

\section{Influence sur la traînée}

Une nutation importante augmente généralement :

\begin{itemize}
    \item l'incidence moyenne ;
    \item la traînée ;
    \item la dispersion.
\end{itemize}

Une réduction rapide de la nutation est donc souhaitable.

\section{Influence sur la précision}

La nutation participe aux écarts observés sur les trajectoires réelles.

Son influence dépend notamment :

\begin{itemize}
    \item de la stabilité gyroscopique ;
    \item de la qualité du projectile ;
    \item des perturbations initiales.
\end{itemize}

\section{Projectile bien stabilisé}

Un projectile correctement stabilisé présente généralement :

\begin{itemize}
    \item une faible nutation ;
    \item un amortissement rapide ;
    \item une bonne précision.
\end{itemize}

\section{Projectile mal stabilisé}

Si la stabilité est insuffisante :

\begin{itemize}
    \item la nutation peut persister ;
    \item l'amplitude peut augmenter ;
    \item la précision se dégrade fortement.
\end{itemize}

Dans les cas extrêmes, le projectile peut devenir instable.

\section{Observation expérimentale}

La nutation est difficile à observer directement à l'œil nu.

Son étude repose souvent sur :

\begin{itemize}
    \item la photographie ultra-rapide ;
    \item les radars Doppler ;
    \item les mesures instrumentées ;
    \item la simulation numérique.
\end{itemize}

\section{Nutation et simulation}

Les modèles simples ignorent généralement la nutation.

Les modèles avancés peuvent la reproduire.

Elle apparaît naturellement dans les modèles :

\begin{itemize}
    \item 5DOF ;
    \item 6DOF.
\end{itemize}

\begin{simulation}

La reproduction correcte de la nutation constitue souvent un indicateur de
qualité pour les modèles balistiques avancés.

\end{simulation}

\section{Vers la dérive gyroscopique}

La rotation, la précession et la nutation constituent les principaux
mouvements dynamiques du projectile.

Ils contribuent également à l'apparition d'un phénomène très connu :

la dérive gyroscopique ou \emph{spin drift}.

Ce phénomène sera étudié dans le chapitre suivant.

\section{À retenir}

\begin{aretenir}

\begin{itemize}
    \item La nutation est une oscillation superposée à la précession.
    \item Elle apparaît dès la sortie du canon.
    \item L'atmosphère tend à amortir progressivement ce mouvement.
    \item Une nutation excessive dégrade la précision.
    \item Les modèles avancés permettent de la reproduire.
\end{itemize}

\end{aretenir}
